\documentclass[10pt,a4paper]{amsart}
\usepackage[margin=19mm]{geometry}
\usepackage{amsmath,amssymb,mathtools}
\usepackage[scr=rsfs,cal=euler]{mathalfa}
\usepackage{xfrac}
\usepackage[unicode,hidelinks,bookmarks=false]{hyperref}
\newcommand{\ie}{i.e.\ }
\newcommand{\cf}{cf.\ }
\newcommand{\resp}{respectively\ }
\newcommand{\Subsection}{\subsection}
\providecommand{\nobreakdash}{\nobreak\hbox{-}}
\setlength{\emergencystretch}{2em}
\multlinegap0pt
\usepackage{bm}
\usepackage{bbm}
\usepackage{dsfont}
\usepackage{xspace}
\usepackage{cancel}
\usepackage{mathtools}
\usepackage{amsthm}
\newtheorem{theorem}{Theorem}
\newtheorem{claim}[theorem]{Claim}
\title[PICPIs: Prediction-Interval-Conditional Prediction Intervals]{PICPIs: Prediction-Interval-Conditional Prediction Intervals}
\author{Xuelin Yang, Baihe Huang, Yilong Hou, Guido Imbens, Michael I. Jordan}
\date{Source: arXiv:2609.25388v1 (CC BY 4.0)}
\begin{document}
\maketitle

\noindent\emph{Source and licence.} PICPIs: Prediction-Interval-Conditional Prediction Intervals. Version arXiv:2609.25388v1 (CC BY 4.0), distributed under the Creative Commons Attribution 4.0 International licence, as recorded in the arXiv licence metadata for this exact version and shown on the abstract page https://arxiv.org/abs/2609.25388v1. Full rights notice: licenses/arxiv-2609.25388-CC-BY-4.0.txt.\par\smallskip

\noindent\textit{Editorial scope.} Counted unit: the Claim whose source label is \texttt{cla:kl-regularity} in the article (source0.tex lines 813--828, source label \texttt{cla:kl-regularity}), delivered together with the article's own complete original proof of it (source0.tex lines 830--878, one \texttt{proof} environment of 49 lines). No mathematics is written, repaired or re-derived: statement and proof are the article's own text line for line. Required context retained uncounted: none. Every definition, display and label of those lines is retained unchanged. Omitted from this package and not redistributed: 1--812, 829--829, 879--3304. Nothing inside a retained excerpt is omitted or altered: every retained body is delivered exactly as the article's lines are written, so no omission receipt of any kind accompanies this package. Citations: the retained excerpts carry no citation command at all, so no bibliography entry of the article is transcribed and no reference is redistributed. Reference adaptation: none was needed and none was made; every reference inside the delivered unit resolves inside it, so no unresolved reference or citation is produced. Printed slips kept literal: no printed word, symbol, hypothesis or number of the counted statement or of its proof was altered. Layout and class: the article's own class is \texttt{article} with packages including geometry, babel, fontenc, newtxtext, microtype, setspace, titlesec, etoolbox, amsmath, amssymb, amsfonts, amsthm, mathtools, newtxmath; the wrapper is the lane's pinned \texttt{amsart} editorial wrapper at 10pt on A4, which restates the theorem-like environments the retained text needs and adds the article's own macro definitions that the retained text needs (none), with extra wrapper packages (bm, bbm, dsfont, xspace, cancel, mathtools). The delivered numbering is the unit's own. Assets: no retained span contains a figure environment, a graphics-inclusion command, a PDF-page inclusion, a table or a TikZ picture; no image file is copied into this package, the package declares no asset dependency and no image file is redistributed with it. Licence: version arXiv:2609.25388v1 of this article is distributed under the Creative Commons Attribution 4.0 International licence, as recorded in the arXiv licence metadata for this exact version and shown on the abstract page https://arxiv.org/abs/2609.25388v1. A full rights notice is delivered at licenses/arxiv-2609.25388-CC-BY-4.0.txt. Characters: every retained excerpt is pure ASCII, so the delivered engine, XeLaTeX with its default Latin Modern text font, reports no missing character.
\begin{claim}[Nearly-regularity]
\label{cla:kl-regularity}
Let $\mu_B$ denote Lebesgue measure on $[0,1]$. Let $X^*$ and $X$ be random variables supported on $[0,1]$ with laws $\nu^*$ and $\nu$, respectively. Assume:
\begin{itemize}
  \item $X^*$ is $\lambda$-regular, i.e. $\nu^*(A) \ge \lambda \mu_B(A)$ for every measurable $A\subseteq[0,1]$.
  \item $\nu^*\ll \nu$ and $D_{\mathrm{KL}}(\nu^*\|\nu)\le \varepsilon$.
\end{itemize}
Then for every $t\in(0,1)$ there exists a measurable set $B_t\subseteq[0,1]$ such that:
\begin{itemize}
  \item $\mu_B(B_t)\le \dfrac{\varepsilon}{\lambda \log(1/t)}$,
  \item for every measurable $A\subseteq [0,1]\setminus B_t$, we have $\nu(A)\ge t\lambda\,\mu_B(A)$.
\end{itemize}
Equivalently, $X$ is $(t\lambda)$-regular on the complement of $B_t$, where the exceptional set has Lebesgue measure at most $\varepsilon/(\lambda\log(1/t))$.

In particular, with $t=\tfrac12$, there exists $B\subseteq[0,1]$ with $\mu_B(B)\le \varepsilon/(\lambda\log 2)$ such that $\nu(A)\ge (\lambda/2)\mu_B(A)$ for all measurable $A\subseteq[0,1]\setminus B$.
\end{claim}

\begin{proof}
Since $X^*$ is $\lambda$-regular, $\nu^*$ dominates $\mu_B$ with constant $\lambda$. In particular, $\nu^*\ll \mu_B$ and the Radon--Nikodym derivative $f^* := \frac{d\nu^*}{d\mu_B}$ exists and satisfies $f^*(x)\ge \lambda$ for $\mu_B$-a.e.\ $x$.

Because $\nu^*\ll \nu$, the Radon--Nikodym derivative $g := \frac{d\nu^*}{d\nu}$ exists $\nu$-a.e.\ and the reverse KL divergence is
\begin{align*}
D_{\mathrm{KL}}(\nu^*\|\nu)
&~:= \int_{[0,1]} \log\!\left(\frac{d\nu^*}{d\nu}\right)\, d\nu^*
 = \int_{[0,1]} \log(g)\, d\nu^*.
\end{align*}
Also $\nu\ll \mu_B$ on the support of $\nu^*$ (and hence on $[0,1]$ after redefining on a null set), so we may write $f:=\frac{d\nu}{d\mu_B}$. Then wherever both densities exist and $f(x)>0$,
\begin{align*}
g(x)=\frac{d\nu^*/d\mu_B}{d\nu/d\mu_B}=\frac{f^*(x)}{f(x)}.
\end{align*}

Fix $t\in(0,1)$ and define the exceptional set
\begin{align*}
B_t := \{x\in[0,1]: f(x) < t\lambda\}.
\end{align*}
On $B_t$, we have $f^*(x)\ge \lambda$ and $f(x)<t\lambda$, hence
\begin{align*}
\frac{f^*(x)}{f(x)} \ge \frac{\lambda}{t\lambda}=\frac1t,
\qquad\text{so}\qquad
\log\!\left(\frac{f^*(x)}{f(x)}\right) \ge \log\!\left(\frac1t\right).
\end{align*}
Using this lower bound inside the reverse KL integral yields
\begin{align*}
\varepsilon
\;\ge\; D_{\mathrm{KL}}(\nu^*\|\nu)
&~= \int_{[0,1]} f^*(x)\log\!\left(\frac{f^*(x)}{f(x)}\right)\, d\mu_B(x)\\
&~\ge \int_{B_t} f^*(x)\log\!\left(\frac{f^*(x)}{f(x)}\right)\, d\mu_B(x)\\
&~\ge \int_{B_t} f^*(x)\log\!\left(\frac1t\right)\, d\mu_B(x)\\
&~= \log\!\left(\frac1t\right)\int_{B_t} f^*(x)\, d\mu_B(x)\\
&~= \log\!\left(\frac1t\right)\,\nu^*(B_t)\\
&~\ge \log\!\left(\frac1t\right)\,\lambda\,\mu_B(B_t),
\end{align*}
where the last step uses $\lambda$-regularity of $\nu^*$ applied to the set $B_t$. Rearranging gives
\begin{align*}
\mu_B(B_t)\le \frac{\varepsilon}{\lambda\log(1/t)}.
\end{align*}

Finally, for any measurable $A\subseteq [0,1]\setminus B_t$, we have $f(x)\ge t\lambda$ for $\mu_B$-a.e.\ $x\in A$, hence
\begin{align*}
\nu(A)
&~= \int_A f(x)\, d\mu_B(x)
\;\ge\; \int_A t\lambda\, d\mu_B(x)
~= t\lambda\,\mu_B(A).
\end{align*}
This proves that $\nu$ is $(t\lambda)$-regular on $[0,1]\setminus B_t$ with the stated bound on $\mu_B(B_t)$.
\end{proof}

\end{document}
